%==================================================================
% Ini adalah lembar daftar penilaian pembimbing lapangan
%==================================================================

%% DILARANG EDIT BAGIAN INI
\newpage
\addcontentsline{toc}{chapter}{LEMBAR DAFTAR PENILAIAN \MakeUppercase{{\tipelaporan} PEMBIMBING LAPANGAN}}
\begin{center}
    Daftar Penilaian {\tipelaporan}\\[0.2cm]
    Pembimbing Lapangan\\
\end{center}

\vspace{\baselineskip}

\begin{flushleft}
    \begin{tabular}{@{}p{3cm}rp{8cm}@{}}
        Nama                        & : & {\penulis} \\
        NIM                         & : & {\nim} \\
        Program                     & : & {\gelarsingkat} - {\prodi} \\
        Nama Instansi               & : & {\namainstansi} \\
        Alamat Instansi             & : & {\alamatinstansi} \\
        Judul \tipelaporansingkatan & : & {\judullaporanid} \\
    \end{tabular}
\end{flushleft}

\begin{center}
    \begin{tabularx}{\textwidth}{|
        m{7mm}|
        m{2.5cm}|
        >{\centering\arraybackslash}X|
        >{\centering\arraybackslash}X|
        >{\centering\arraybackslash}X|
        m{11mm}|
        m{11mm}|}
        \hline
        \multirow{2}{*}{\centering No} &
        \multirow{2}{*}{\centering Uraian} &
        \centering Baik Sekali &
        \centering Baik &
        \centering Cukup &
        \multirow{2}{*}{\centering Bobot} &
        \multirow{2}{*}{\centering Nilai} \\
        \cline{3-5}
        & & (80–100) & (66–79.99) & (60–65.99) & & \\ 
        \hline

        \rule{0pt}{3ex}
        \centering 1 &
        \centering Kehadiran Mahasiswa / Konsultasi &&
        &&
        \centering 10\% &
        \\ \hline

        \rule{0pt}{3ex}
        \centering 2 &
        \centering Kemandirian &&
        &&
        \centering 30\% &
        \\ \hline

        \rule{0pt}{3ex}
        \centering 4 &
        \centering Pemahaman Materi \tipelaporansingkatan &&
        &&
        \centering 40\% &
        \\ \hline

        \rule{0pt}{3ex}
        \centering 3 &
        \centering Sistematika Penulisan Laporan &&
        &&
        \centering 20\% &
        \\ \hline

        \multicolumn{6}{|c|}{\rule{0pt}{3ex}Jumlah Grade} & \\ \hline
    \end{tabularx}
\end{center}

\vspace{\baselineskip}

\begin{flushleft}
    {\kota}, \textbf{\tglpernyataan}\\[0.2cm]
    Dosen Pembimbing\\[3cm]
    ({\dosenpembimbing})
\end{flushleft}
%% DILARANG EDIT BAGIAN INI